
\documentclass[10pt]{article}
\usepackage[utf8]{inputenc}
\usepackage[T1]{fontenc}
\usepackage{amsmath, amssymb, xcolor, geometry, enumitem, multicol, tcolorbox}

\definecolor{mainblue}{RGB}{0, 102, 204}
\geometry{a4paper, margin=1.2cm, top=1cm, bottom=3cm, footskip=1.2cm}

\begin{document}
\enlargethispage{2.5cm}

% Modern Header
\begin{tcolorbox}[colback=mainblue!5, colframe=mainblue, sharp corners, boxrule=0.5pt]
    \small \textbf{Unit:} undefined \hfill \textbf{Name:} \rule{3cm}{0.4pt} \\
    \small \textbf{Topic:} Variables and Expressions / Order of Operations \hfill \textbf{Date:} \rule{2cm}{0.4pt} \\
    \centering \textbf{\large Challenge (Level 0)}
\end{tcolorbox}

\vspace{0.2cm}

% MCQ Section
\begin{multicols}{2}
\begin{enumerate}[label=\textbf{Q\arabic*.}, leftmargin=*, itemsep=6pt, parsep=0pt]

  \item \small What is the variable in the expression $2x + 3$?
  \begin{enumerate}[label=(\Alph*), leftmargin=0.6cm, nosep, itemsep=2pt]
    \item 2
    \item 3
    \item x
    \item 2x
    \item 2x + 3
  \end{enumerate}

  \item \small Identify the coefficient in the term $4y$?
  \begin{enumerate}[label=(\Alph*), leftmargin=0.6cm, nosep, itemsep=2pt]
    \item y
    \item 4
    \item 4y
    \item 1
    \item None
  \end{enumerate}

  \item \small What is the constant term in the expression $x - 2$?
  \begin{enumerate}[label=(\Alph*), leftmargin=0.6cm, nosep, itemsep=2pt]
    \item x
    \item -2
    \item 2
    \item 1
    \item 0
  \end{enumerate}

  \item \small Identify the operator in $7 	imes 2$?
  \begin{enumerate}[label=(\Alph*), leftmargin=0.6cm, nosep, itemsep=2pt]
    \item 7
    \item 2
    \item $	imes$
    \item 14
    \item None
  \end{enumerate}

  \item \small What is the expression $3 + 2$ an example of?
  \begin{enumerate}[label=(\Alph*), leftmargin=0.6cm, nosep, itemsep=2pt]
    \item A variable
    \item A constant
    \item An operation
    \item A simple expression
    \item An equation
  \end{enumerate}

  \item \small Identify the operation being performed in $9 - 1$?
  \begin{enumerate}[label=(\Alph*), leftmargin=0.6cm, nosep, itemsep=2pt]
    \item Addition
    \item Subtraction
    \item Multiplication
    \item Division
    \item Exponentiation
  \end{enumerate}

  \item \small What does the term $2x^2$ represent?
  \begin{enumerate}[label=(\Alph*), leftmargin=0.6cm, nosep, itemsep=2pt]
    \item A constant
    \item A variable
    \item An algebraic expression
    \item A geometric term
    \item An exponential term
  \end{enumerate}

  \item \small In the expression $5(2x)$, what operation is being performed between $5$ and $2x$?
  \begin{enumerate}[label=(\Alph*), leftmargin=0.6cm, nosep, itemsep=2pt]
    \item Addition
    \item Subtraction
    \item Multiplication
    \item Division
    \item Exponentiation
  \end{enumerate}

\end{enumerate}
\end{multicols}

\vspace{-0.2cm}
\hrule
\vspace{0.2cm}
\textbf{\small Open Ended Questions (FRQ)}
\begin{enumerate}[label=\textbf{Q\arabic*.}, start=9, itemsep=5pt, topsep=0pt]

  \item \small Draw a simple diagram to illustrate the concept of a variable and explain its role in an algebraic expression. \vspace{2cm}

  \item \small Explain the difference between a constant and a variable in the context of algebraic expressions. Provide at least two examples. \vspace{2cm}

\end{enumerate}

\vfill

% Chef's Tips Footer
\begin{tcolorbox}[colback=blue!5, colframe=mainblue!50, title=\small \textbf{Chef's Tips}, fonttitle=\bfseries, bottom=1mm]
\begin{itemize}[noitemsep, topsep=2pt, leftmargin=0.5cm]
  \item \scriptsize Emphasize the distinction between variables and constants in algebraic expressions.\item \scriptsize Ensure students understand the role of coefficients in terms with variables.\item \scriptsize Use simple examples to illustrate the order of operations and how it applies to basic expressions.
\end{itemize}
\end{tcolorbox}

\end{document}
  